% !TEX root = ../main.tex
\section{Descrizione del tema e del contesto}
\label{sec:contesto}

\subsection{Dipendenze e confini di fiducia}
\subsubsection{Dipendenze dirette e transitive in npm}

\textbf{npm} è il gestore di pacchetti più usato nell'ecosistema JavaScript.
Un pacchetto è un componente software pubblicato da un manutentore e riutilizzato
da altri sviluppatori. Una dipendenza è \emph{diretta} quando il progetto la
dichiara e installa esplicitamente; è \emph{transitiva} quando è richiesta da
un'altra dipendenza. Chi sviluppa un'applicazione si affida quindi sia ai
manutentori dei pacchetti scelti sia alle dipendenze che questi includono.
L'account con cui il manutentore pubblica le versioni è un punto sensibile:
chi ne ottiene il controllo può distribuire aggiornamenti agli utilizzatori.

\subsubsection{Ricette di build e installazione in AUR}
L'Arch User Repository (AUR) ospita ricette comunitarie, in particolare file
\texttt{PKGBUILD}, non i pacchetti binari dei repository ufficiali.
\texttt{makepkg} usa la ricetta per recuperare, compilare e impacchettare il
programma; prima della costruzione l'utente deve esaminare il file e gli
eventuali script associati~\cite{archwiki-aur,archwiki-makepkg}. In npm si
integra una libreria nel progetto; con AUR si seguono istruzioni per creare e
installare un pacchetto sul sistema.

\subsection{Soggetti e asset interessati}
\subsubsection{Chi può essere esposto}

I soggetti rilevanti sono il manutentore, il cui account o pacchetto può essere
preso in controllo; lo sviluppatore, che integra dipendenze nel proprio
progetto; e l'utente finale o di sistema, sul cui browser o computer può essere
eseguito il codice distribuito.

\subsubsection{Che cosa occorre proteggere}
Gli asset rilevanti sono account di pubblicazione, integrità di pacchetti e
ricette, ambienti di build e installazione e dati locali accessibili al
processo. Il rischio varia in base al contesto e ai privilegi d'esecuzione.
